\documentclass{article}
\usepackage{amssymb}
\usepackage{amsfonts}
\usepackage{amsmath}

\author{Jonathan Kalsky}
\date{March 5 2026}
\title{\normalsize CSCE 222\\ \Large Homework 3}
\linespread{1.5}

\begin{document}
\maketitle
\tableofcontents
\vspace{1cm}
\section{Section 2.3}
\subsection{Question 12}
Determine whether each of these functions from $\mathbb{Z}$ to $\mathbb{Z}$
is one-to-one.
\subsubsection{part a}
$f(n) = n - 1$\\
Let $a,b \in \mathbb{Z} \mid a \neq b$\\
$(a \neq b) \Rightarrow (a-1 \neq b-1)$ under subtraction\\
This concludes that 
$(a \neq b) \iff f(a) \neq f(b)$\\
$\therefore$ the function is one to one by definition

\subsubsection{part b}
$f(n) = n^2 + 1$\\
Let $a \in \mathbb{Z}$\\
$\exists b \in \mathbb{Z} | b = (-1)a \neq a$ | Integers are closed under scalar multiplication\\
$f(a) = a^2 + 1 = (-a)^2 + 1 = b^2 + 1 = f(b)$ | Substituion\\
$\exists a,b ((a \neq b) \land (f(a) = f(b)))$
$\therefore$ the function is not one to one by definition

\subsubsection{part c}
$f(n) = n^3$\\
Let $a,b \in \mathbb{Z} \mid a \neq b$\\
$f(a) = a^3$ and $f(b) = b^3$ | Substitution\\
$\sqrt[3]{a^3} = a \neq b = \sqrt[3]{b^3}$ | the cube root has a unique solution\\
$\therefore \forall a,b \in \mathbb{Z}, (f(a) \neq f(b) \iff a \neq b)$  so $n^3$ is injective by definition

\subsubsection{part d}
$f(n) = \lceil n/2 \rceil $\\
let a = 3, and b = 4\\
$f(a) = f(3) = 2$ and $f(b) = f(4) = 2$\\
$\exists a,b \in \mathbb{Z}, (a \neq b) \land (f(a) \neq f(b))$
$\therefore$ the function is not one to one by definition

\subsection{Question 14 Determine whether f : Z ×Z → Z is onto}
\subsubsection{part a}
$f(m,n) = 2m - n$\\
Let m = 0\\
f(0,n) = -n\\
since $n \in \mathbb{Z}$, f(0,n) spans all $\mathbb{Z}$
$\therefore$ The function is onto

\subsubsection{part b}
$f (m, n)= m^2 - n^2$\\
$f (m, n)= (m-n)(m+n)$\\
$b \in \mathbb{Z}$\\
Let b = 2 = (a)(b) $\mid a,b \in \mathbb{Z}$\\
2 = (2)(1) $\lor$ (1)(2)\\ 
((m-n = 2) $\land$ (m+n = 1)) $\lor$ ((m+n = 2) $\land$ (m-n = 1))\\
Assuming (m-n = 2), m = 2+n, 2n+2 = 1 $\to n \notin \mathbb{Z}$\\
So (m-n = 2) $\land$ (m+n = 1) is False\\
Assuming (m+n = 2), m = 2 - n, 2-2n = 1 $\to n \notin \mathbb{Z}$\\
So (m+n = 2) $\land$ (m-n = 1) is False as well\\
Since there is an integer, that does not have a preimage in the domain, this is not an onto function
\subsubsection{part c}
\subsubsection{part d}
\subsubsection{part e}
\subsection{Question 20 Give an example of a function from N to N that is}
\subsubsection{part a. one-to-one but not onto.}
$f(x) = x^2$
\subsubsection{part b. onto but not one-to-one.}
$f(x) = \lceil x/2 \rceil$
\subsubsection{part c. one-to-one and onto, but not the identity.}
$f(x) = x + (-1)^{x+1}$
\subsubsection{part d. neither one-to-one nor onto}
$f(x) = sin(x)$
\subsection{Question 22 Determine whether each of these functions is a bijection
from R to R}
\subsubsection{part a. $f (x)= -3x +4$}
Let y = -3x + 4\\
$f^{-1}(y) = \frac{y-4}{-3}$\\
Since we are able to determine a unique inverse the function is a bijection
\subsubsection{part b. $f (x)= -3x^2 +7$}
\subsubsection{part c. $f (x)= (x +1)/(x +2)$}
Assuming f is bijective, $f^{-1}(y) = \frac{2y}{1-y}$\\
This shows y is not defined in the range of f(x), causing division by 0.\\
By contradiction, the range is not equal to the codomain, and thus by failing surjection, the function is not a bijection.
\subsubsection{part d}
\subsection{Question 60. How many bytes are required to encode n bits of data
where n equals}
\subsubsection{part a. n = 4}
$4/8 = 0.5$, thus only 1 byte is required
\subsubsection{part b. n = 10}
$10/8 = 1.25$, thus 2 bytes are required
\subsubsection{part c. n = 500}
$500/8 = 62.5$, thus 63 bytes are required
\subsubsection{part d. n = 3000}
$3000/8 = 300+70+5 = 375$, thus 375 bytes are required


\end{document}